\section{Discussion}
\label{sec:discussion}

There are several recent examples of supply chain vulnerabilities due to shortages in raw and intermediate materials. \tool{}s
are an important means of discovering and exploring such systemic effects.
For example, the COVID-19 pandemic exposed fundamental vulnerabilities in semiconductor supply chains, where hidden intermediates and concentrated manufacturing led to cascading global shortages. Hasan et al.~\cite{hasan2022demand} documented how automotive manufacturers faced severe disruptions because chip suppliers had reallocated production capacity to consumer electronics, revealing how obscured intermediate dependencies amplified supply chain vulnerabilities across multiple industries. The semiconductor shortage demonstrated that traditional top-down visibility failed to capture the complex web of intermediate processing steps and their interdependencies, leading to unexpected bottlenecks~\cite{chakraborty2022empirical,tam2022understanding}.

Similarly, China's dominance in rare earth elements processing illustrates how intermediate transformation stages can create hidden supply chain risks. Zhang et al.~\cite{zhang2024reshaping} found that terbium production in China declined by 90\% between 2007 and 2018, not due to resource scarcity but because existing mining techniques could not meet stringent environmental regulations. This research revealed that only 25\% of China's rare earth production quota was utilized by 2018, indicating that environmental constraints at intermediate processing stages—rather than political restrictions—had become the primary supply bottleneck. The study highlighted how top-down approaches focusing on end products missed critical vulnerabilities in upstream processing capabilities.

The lithium battery supply chain presents another compelling example of hidden intermediate dependencies. Hu et al.~\cite{hu2025technological} demonstrated that technological risks and export controls on advanced battery manufacturing processes can disrupt entire supply networks, with the United States holding 56-79\% of core technological capabilities. Their analysis revealed that restrictions on cathode formulation technologies and intellectual property transfers created bottlenecks affecting downstream battery production across multiple countries, despite abundant lithium ore availability. The lithium battery case shows that proprietary manufacturing processes and technology dependencies create supply chain vulnerabilities that are invisible to traditional backward-tracing approaches.

The disruptions to semiconductor, rare earth, and lithium batteries demonstrate a fundamental limitation of top-down supply chain analysis: the methodology inherently struggles with opacity at multiple levels of the transformation network. Baldwin et al.~\cite{baldwin2023hidden} provided empirical evidence that traditional face-value supply chain analysis significantly underestimates exposure to foreign suppliers due to hidden intermediates, finding that the U.S. vehicle sector's exposure to Chinese industrial inputs was four times higher than indicated by conventional top-down measures. This ``hidden exposure'' stems from the complex spiral of inputs used to make inputs, which remains largely invisible when starting from finished products and working backward through proprietary supplier networks.